\section{Localization under support restriction}
\label{sec:support-restriction}

We now compare localization claims made on nested regimes while holding the network, receiver decomposition, and ambient phenomenon fixed.

Fix an ambient source regime $C_\ast\subseteq X_0$ and a phenomenon profile
\[
\Phi_r^{\ast}:S_r^{\theta,C_\ast}\to Y_\Phi.
\]
For each nonempty $C\subseteq C_\ast$, restrict $\Phi_r^{\ast}$ to $S_r^{\theta,C}$ and pull it back to earlier cuts. We write the resulting profile as $\Phi_q^{\theta,C}$.

If $C\subseteq C'\subseteq C_\ast$, then
\[
S_q^{\theta,C}\subseteq S_q^{\theta,C'}
\]
and both the receiver and phenomenon equivalence relations restrict to the smaller domain:
\begin{align}
D_{q,j}^{\theta,C}
&=
D_{q,j}^{\theta,C'}
\cap
\bigl(S_q^{\theta,C}\times S_q^{\theta,C}\bigr),
\label{eq:receiver-kernel-restriction}\\
D_{\Phi,q}^{\theta,C}
&=
D_{\Phi,q}^{\theta,C'}
\cap
\bigl(S_q^{\theta,C}\times S_q^{\theta,C}\bigr).
\label{eq:phenomenon-kernel-restriction}
\end{align}
The network has not acquired a new factorization. The criterion is being tested on fewer state pairs.

\begin{theorem}[Realization is preserved by support restriction]
\label{thm:realization-monotonicity}
Let $C\subseteq C'\subseteq C_\ast$. If a receiver family $K$ realizes $\Phi$ on $C'$, then $K$ realizes $\Phi$ on $C$. Equivalently,
\begin{equation}
\boxed{
\Rcal_{\Phi,q}^{\theta,C'}
\subseteq
\Rcal_{\Phi,q}^{\theta,C}.
}
\label{eq:realization-family-monotonicity}
\end{equation}
\end{theorem}

\begin{proof}
Assume $K\in\Rcal_{\Phi,q}^{\theta,C'}$. Restrict
\[
\bigcap_{j\in K}D_{q,j}^{\theta,C'}
\subseteq
D_{\Phi,q}^{\theta,C'}
\]
to $S_q^{\theta,C}\times S_q^{\theta,C}$ and use \eqref{eq:receiver-kernel-restriction}--\eqref{eq:phenomenon-kernel-restriction}.
\end{proof}

\begin{corollary}[Narrower support cannot increase the localization order]
\label{cor:order-monotonicity}
Under the hypotheses of \cref{thm:realization-monotonicity},
\begin{equation}
\boxed{
 d_{\Phi,q}^{\theta,C}
\le
 d_{\Phi,q}^{\theta,C'}.
}
\label{eq:order-monotonicity}
\end{equation}
\end{corollary}

The theorem gives a one-sided comparison. A broad regime demands that the proposed receiver family preserve the phenomenon across more possible internal distinctions. Restricting the regime removes some of those comparisons, so a family that failed globally may become sufficient locally. The reverse inference is invalid without additional information.

The continuous example in \cref{sec:supported-process} is already strict. For the linear readout $\Phi(x_1,x_2)=x_1+x_2$ with coordinate accesses, the full domain $\mathbb R^2$ requires both coordinates, while the diagonal $D=\{(t,t)\}$ has localization order one. The phenomenon remains nonconstant. The decrease is caused entirely by the relation $x_1=x_2$ on the restricted support.

\subsection{Minimal identity can change without changing the count}

Support also changes which receiver families are minimal. Consider the ambient phenomenon
\[
\Phi(x_1,x_2)=x_1
\]
on $\mathbb R^2$. The first coordinate is the unique one-coordinate realization on the full domain. On the diagonal $D$, either coordinate realizes the same phenomenon because $x_1=x_2$. The minimum count remains one, but the set of minimal realizations changes from
\[
\bigl\{\{1\}\bigr\}
\quad\text{to}\quad
\bigl\{\{1\},\{2\}\bigr\}.
\]
Circuit size therefore carries less information than circuit identity. Agreement in minimum size does not imply agreement in the components that realize the phenomenon, and disagreement in identity does not by itself imply that the ambient computation changed.

At the other extreme, a dependency can disappear completely when the phenomenon becomes constant on a restricted support. Taking parity itself as the phenomenon, the example from \cref{sec:support-dependency} gives such a case. These are exact changes in the feasible localization set, not failures of an extraction algorithm.

\subsection{Natural and intervention supports}

\Cref{thm:realization-monotonicity} is stated for source-generated regimes because those supports come with compatible continuations across cuts. At a fixed cut, its proof uses only restriction of the receiver and phenomenon relations. For arbitrary analysis supports $S\subseteq S'\subseteq X_q$ in the sense of \cref{def:analysis-support}, the same implication holds whenever the same ambient access and phenomenon are restricted to both domains.

Intervention protocols make this distinction concrete. Let $S_q^{\mathrm{nat}}$ be an analysis support generated by natural inputs, and let $S_q^{\mathrm{int}}$ contain states admitted by a patching, ablation, or other intervention protocol. The second set need not be the prefix image of any source regime. The statements
\[
Q_{q,j}^\theta|_{S_q^{\mathrm{nat}}}
\quad\text{realizes }\Phi|_{S_q^{\mathrm{nat}}}
\]
and
\[
Q_{q,j}^\theta|_{S_q^{\mathrm{int}}}
\quad\text{realizes }\Phi|_{S_q^{\mathrm{int}}}
\]
are therefore claims on different restricted domains.

If one support contains the other, exact realization transports from the larger domain to the smaller one. If the domains are incomparable, support restriction alone gives no implication in either direction. Intervention-induced hidden states can leave the natural representation distribution and can recruit pathways not active on ordinary inputs \citep{makelov2024subspace,grant2026divergent}. In the present formalism, those intervention states belong to the domain on which the localization claim is stated.
